\documentclass[11pt]{article}
\usepackage[margin=2.2cm]{geometry}
\usepackage{booktabs,graphicx,float,textcomp,longtable,hyperref}
\title{Context composition for intrusion detection with scarce attack labels}
\author{Research record: intent, selected evidence, and evaluation design}
\date{2 October 2026; revision r06}
\begin{document}
\maketitle
\section{Research intent and claim to test}
The primary objective is to improve detection of attacks for which only a small number of labelled examples are available. Improving every class or maximising aggregate macro-F1 is not the central claim. Under a limited context budget, the research question is which common and attack-specific examples help discriminate these attacks from benign traffic and other attacks while preserving existing detection performance.

TabPFN is a candidate because it conditions predictions on labelled context examples using a pretrained backbone. This permits context changes without updating backbone weights. It does not establish better few-label performance, lower total cost, or immediate online adaptation. XGBoost remains a strong baseline and, when model updates are studied, can use continued training as well as refitting. Preprocessing, context construction, caching, inference, calibration and any learned context or embedding must be charged to the relevant method.

\paragraph{Working hypothesis.} A context composition strategy can make scarce attack labels more useful than an unstructured context with the same information and memory budget. The resulting class-level benefit must be measured together with benign false alarms, other-class performance and computational cost. Scarcity is defined from available training labels, before examining the final evaluation outcomes.

\paragraph{Candidate mechanisms.} Diverse benign examples may better represent normal behaviour. Tail synthesis may broaden scarce attack coverage, with real-example duplication as a necessary control. Learned representations may improve class separation. Context optimisation may allocate limited context capacity more effectively. These are separate hypotheses, not established findings or a commitment to combine every mechanism. Relevant starting points include \href{https://proceedings.iclr.cc/paper_files/paper/2024/file/5d54d2df6ec8f7b920aa0fec9a6d1b2e-Paper-Conference.pdf}{LITO (ICLR 2024)}, \href{https://github.com/mlopezm/Supervised-contrastive-learning-over-prototype-label-embeddings}{supervised contrastive prototype embeddings}, and \href{https://arxiv.org/abs/2402.06971}{in-context data distillation}. The latter is associated with the ICLR 2024 ME-FoMo workshop, not the main conference.

\section{Evaluation direction: sequential attack introduction}
\textbf{Design status: not yet an executed or fixed protocol.} The user clarified that reproducing raw collection timestamps is not mandatory, and that comparing performance immediately before and after acquiring a few labels is not a required objective. Sequential introduction is an evaluation axis for scarce-label attack detection; online context updating is a separate design choice.

A conceptual evaluation starts with benign traffic and some attacks, introduces additional attacks in stages, and observes class-specific performance as the mixture changes. An attack first appearing in the evaluation stream may already have a few training examples. It must not be called unseen or out-of-distribution solely because it appears later in the test sequence. Attack order, prevalence, available training labels and context update policy must be explicitly documented and controlled. With a fixed context and the same samples, reordering evaluation alone leaves individual predictions unchanged. Such a scenario measures performance under stage-specific mixtures; it is not evidence of adaptation without an update mechanism. No fixed number of stages, label quota or update interval is adopted in this revision.

\paragraph{Comparisons.} Use the same labelled information and evaluation rows for XGBoost, Global TabPFN, and the proposed context composition. Match total stored context, including repeated anchors, when isolating composition or expert-count effects. Report additional data used to train embeddings, optimise contexts or fit scorer/verifier policies. Introduce component ablations after a useful configuration is identified on validation data.

\paragraph{Measurements.} Report per-attack precision, recall and F1, benign false alarms, retained performance on existing attacks, and context construction/model fitting/inference time and memory. Benign false-alarm rate is the fraction of true benign examples predicted as attacks; it is not the fraction of attacks predicted benign. Report support and prevalence in each stage. Stage macro-F1 is secondary and cannot be compared without accounting for differing class sets. Avoid selecting a preferred attack order or model setting on the final test outcomes.

\paragraph{Optional updates and OOD extension.} If online updates are later introduced, prediction must precede label availability, and only labels available at that moment may affect future predictions. A globally chronological split is needed when claiming natural-time deployment performance. Later OOD evaluation should separately address distribution shifts or attacks absent from training, distinguishing recognition before labels from adaptation after labels. Few-label classification, sequential arrival and zero-label novelty detection are distinct problems.

\section{Existing exploratory evidence}
All reported results below use seed 43 and the conflict-free development holdouts already inspected during model development. They are not final, untouched evaluations. Removing feature-identical rows with conflicting labels reduced the full CIC-IDS2018 dataset from 20,115,529 to 14,755,417 rows, and ToN-IoT from 27,520,260 to 10,807,288 rows. These full-dataset counts are distinct from the test supports in the tables. Feature-label conflicts motivate auditing; they do not establish that a remaining class is intrinsically impossible to classify.

The implemented context design combines a common multi-class anchor with a residual-cluster-specific block. Global errors help construct the experts offline. Scorer and verifier select whether to call an expert and accept its output. Their policies were selected on validation data. Changing the context can require revalidating these policies; adding examples does not automatically validate routing.

The current evidence supports a narrower observation: ToN Scanning improves under the residual-context system, while the selected CIC-IDS2018 policies retain Global predictions. It does not establish the cause of each tail error or the effectiveness of sequential evaluation. The previous label-assisted assignment diagnostic is not a deployable oracle or a guaranteed performance upper bound. Different assignment rules must not be interpreted as an isolated measure of expert competence.
\begin{table}[H]
\centering
\small
\caption{Macro-F1 and accuracy on the conflict-free test splits (seed 43). The first five methods use exactly the same 100,000 training IDs as the Global context C0; full-train XGBoost and the proposed model use additional training data and are not same-budget comparisons.}
\label{tab:sota_overall}
\resizebox{\linewidth}{!}{%
\begin{tabular}{lrrrrl}
\toprule
Method & CIC macro-F1 & CIC acc. (\%) & ToN macro-F1 & ToN acc. (\%) & Training rows (CIC / ToN) \\
\midrule
Global TabPFN-v3 (C0) & 0.7824 & 99.37 & 0.6796 & 90.13 & 100,000 \\
XGBoost (same 100k) & 0.7821 & 98.83 & 0.6983 & 92.09 & 100,000 \\
BoostPFN & 0.7148 & 98.90 & 0.5866 & 82.68 & 100,000 \\
LoCalPFN (FT) & 0.7024 & 97.72 & 0.6462 & 89.44 & 100,000 \\
DistPFN (v3) & 0.7839 & 99.43 & 0.6682 & 88.56 & 100,000 \\
XGBoost (full train) & 0.7795 & 98.02 & 0.7162 & 92.38 & 8,666,430 / 4,669,119 \\
Proposed: experts (K=4) + S/V & 0.7824 & 99.37 & 0.7617 & 90.31 & 100,000 + expert/route pools \\
\bottomrule
\end{tabular}}
\par\vspace{2pt}\parbox{\linewidth}{\footnotesize Test rows: CIC-IDS2018 3,042,473; ToN-IoT 2,271,723. Source: EXP61 (SOTA), EXP63 K sweep (proposed, K=4: previously used representative configuration, not selected as the test maximum, single seed).}
\end{table}

\begin{table}[H]
\centering
\small
\caption{Per-class F1 on CIC-IDS2018 (seed 43, test rows in descending order). Proposed = experts (K=4) + scorer/verifier (S1V0 policy), EXP63 K sweep.}
\label{tab:classwise_cic2018}
\resizebox{\linewidth}{!}{%
\begin{tabular}{lrrrrrrrr}
\toprule
Class & Test rows & Global TabPFN-v3 & XGBoost (100k) & BoostPFN & LoCalPFN & DistPFN & XGBoost (full) & Proposed K=4 \\
\midrule
Benign & 2,652,203 & 0.9965 & 0.9933 & 0.9937 & 0.9921 & 0.9968 & 0.9887 & 0.9965 \\
DDoS & 262,680 & 0.9922 & 0.9998 & 0.9906 & 0.9433 & 0.9931 & 0.9036 & 0.9922 \\
Bot & 41,542 & 0.9985 & 0.9976 & 0.9942 & 0.9972 & 0.9982 & 0.9977 & 0.9985 \\
DoS & 39,483 & 0.9920 & 0.9998 & 0.9598 & 0.7309 & 0.9919 & 0.9999 & 0.9920 \\
Brute force & 37,696 & 0.9997 & 0.9995 & 0.9997 & 0.9999 & 0.9997 & 0.9996 & 0.9997 \\
Infiltration & 8,742 & 0.2879 & 0.1528 & 0.0653 & 0.1350 & 0.3055 & 0.2927 & 0.2880 \\
Web attack & 127 & 0.2102 & 0.3321 & 0.0000 & 0.1183 & 0.2024 & 0.2745 & 0.2102 \\
\bottomrule
\end{tabular}}
\end{table}

\begin{table}[H]
\centering
\small
\caption{Per-class F1 on ToN-IoT (seed 43, test rows in descending order). Proposed = experts (K=4) + scorer/verifier (S1V0 policy), EXP63 K sweep.}
\label{tab:classwise_toniot}
\resizebox{\linewidth}{!}{%
\begin{tabular}{lrrrrrrrr}
\toprule
Class & Test rows & Global TabPFN-v3 & XGBoost (100k) & BoostPFN & LoCalPFN & DistPFN & XGBoost (full) & Proposed K=4 \\
\midrule
Benign & 1,474,302 & 0.9217 & 0.9385 & 0.9042 & 0.9184 & 0.9064 & 0.9389 & 0.9231 \\
XSS & 343,854 & 0.8305 & 0.8734 & 0.7929 & 0.8469 & 0.8247 & 0.8714 & 0.8305 \\
Password & 210,951 & 0.8921 & 0.8920 & 0.6728 & 0.8911 & 0.8909 & 0.8918 & 0.8921 \\
DDoS & 103,338 & 0.9804 & 0.9842 & 0.9248 & 0.9786 & 0.9850 & 0.9928 & 0.9804 \\
Injection & 68,891 & 0.9027 & 0.9128 & 0.4793 & 0.8925 & 0.9164 & 0.9759 & 0.9047 \\
Backdoor & 40,030 & 0.9990 & 0.9990 & 0.9609 & 0.9988 & 0.9990 & 0.9924 & 0.9990 \\
DoS & 22,458 & 0.9771 & 0.9563 & 0.3954 & 0.5554 & 0.9823 & 0.9814 & 0.9771 \\
Scanning & 6,029 & 0.0394 & 0.0458 & 0.3281 & 0.0973 & 0.0060 & 0.0396 & 0.8587 \\
MitM & 1,161 & 0.1322 & 0.2363 & 0.0468 & 0.1655 & 0.0512 & 0.3236 & 0.1322 \\
Ransomware & 709 & 0.1208 & 0.1449 & 0.3608 & 0.1170 & 0.1204 & 0.1539 & 0.1196 \\
\bottomrule
\end{tabular}}
\end{table}

\begin{table}[H]
\centering
\small
\caption{Measured cost on one RTX 4090 (24\,GB) with up to two jobs in parallel: elapsed times include resource contention and are not standalone latencies. LoCalPFN fit includes validation; its prediction includes kNN retrieval. DistPFN shares the Global fit and prediction.}
\label{tab:sota_cost}
\resizebox{\linewidth}{!}{%
\begin{tabular}{lrrrrc}
\toprule
Method & CIC fit (s) & CIC predict (s) & ToN fit (s) & ToN predict (s) & Peak GPU GiB (CIC / ToN) \\
\midrule
Global TabPFN-v3 (C0) & 36.4 & 1,099.3 & 36.7 & 821.7 & 2.92 / 2.94 \\
XGBoost (same 100k) & 2.3 & 1.0 & 55.4 & 1.0 & 0.55 / 0.53 \\
BoostPFN & 98.5 & 2,308.5 & 97.8 & 1,737.6 & 0.89 / 0.89 \\
LoCalPFN (FT) & 3,196.3 & 17,300.8 & 5,783.1 & 11,144.8 & 7.54 / 7.54 \\
DistPFN (v3) & 36.4 & 1,099.4 & 36.7 & 821.9 & 2.92 / 2.94 \\
XGBoost (full train) & 107.8 & 1.0 & 143.2 & 1.1 & 2.20 / 1.59 \\
\bottomrule
\end{tabular}}
\end{table}

\clearpage
\section{Expert and routing diagnostics}
Experts are compared on identical complete test populations, including false positives from other classes. A target-class improvement should be checked on validation data and against damage to other classes. The K sweep changes total context size and repeats the shared anchor, so it does not isolate the effect of expert count at fixed memory.
\begin{table}[H]
\centering
\small
\caption{Direct expert evaluation on CIC-IDS2018, K=4. All experts predict the same complete test set; class F1 includes false positives from all other classes. This is not label-assisted assignment or an online routing result.}
\label{tab:experts_cic2018}
\begin{tabular}{lrrrrrr}
\toprule
Class & Test rows & Global & Expert 1 & Expert 2 & Expert 3 & Expert 4 \\
\midrule
Benign & 2,652,203 & 0.9965 & 0.9976 & 0.8936 & 0.7945 & 0.8750 \\
DDoS & 262,680 & 0.9922 & 0.9918 & 1.0000 & 0.9999 & 0.9998 \\
Bot & 41,542 & 0.9985 & 0.9985 & 0.9934 & 0.9982 & 0.9980 \\
DoS & 39,483 & 0.9920 & 0.9921 & 0.9861 & 0.9959 & 0.9920 \\
Brute force & 37,696 & 0.9997 & 0.9997 & 0.9997 & 0.9997 & 0.9997 \\
Infiltration & 8,742 & 0.2880 & 0.2290 & 0.0266 & 0.0178 & 0.0227 \\
Web attack & 127 & 0.2102 & 0.3526 & 0.0691 & 0.0770 & 0.1986 \\
\bottomrule
\end{tabular}
\end{table}

\begin{table}[H]
\centering
\small
\caption{Direct expert evaluation on ToN-IoT, K=4. All experts predict the same complete test set; class F1 includes false positives from all other classes. This is not label-assisted assignment or an online routing result.}
\label{tab:experts_toniot}
\begin{tabular}{lrrrrrr}
\toprule
Class & Test rows & Global & Expert 1 & Expert 2 & Expert 3 & Expert 4 \\
\midrule
Benign & 1,474,302 & 0.9217 & 0.9117 & 0.9088 & 0.8580 & 0.9197 \\
XSS & 343,854 & 0.8305 & 0.8248 & 0.8500 & 0.8339 & 0.8497 \\
Password & 210,951 & 0.8921 & 0.8898 & 0.8921 & 0.8907 & 0.8906 \\
DDoS & 103,338 & 0.9804 & 0.9747 & 0.9819 & 0.9682 & 0.9709 \\
Injection & 68,891 & 0.9027 & 0.8505 & 0.9354 & 0.4595 & 0.8730 \\
Backdoor & 40,030 & 0.9990 & 0.9982 & 0.9991 & 0.9970 & 0.9986 \\
DoS & 22,458 & 0.9771 & 0.9624 & 0.9544 & 0.7146 & 0.9687 \\
Scanning & 6,029 & 0.0394 & 0.6157 & 0.5638 & 0.7609 & 0.6605 \\
MitM & 1,161 & 0.1322 & 0.0318 & 0.0389 & 0.0364 & 0.0522 \\
Ransomware & 709 & 0.1208 & 0.0944 & 0.1059 & 0.1322 & 0.1166 \\
\bottomrule
\end{tabular}
\end{table}

\begin{table}[H]
\centering
\small
\caption{Scorer/verifier (S1V0) on the seed-43 banks with K = 2, 4, 6, 8 (EXP62/63). Calls and acceptances are logical counts obtained by applying the validation-selected policy to stored expert predictions. CIC-IDS2018 keeps the Global prediction at every K; ToN-IoT calls one expert per row ; Scanning is a prominent improvement. Total context rows differ across K, so this is not a fixed-memory comparison.}
\label{tab:sv_by_k}
\begin{tabular}{llrrrrrr}
\toprule
Dataset & Bank & Macro-F1 & Call rate (\%) & Accept rate (\%) & Changed & Corrected & Damaged \\
\midrule
CIC-IDS2018 & Global & 0.7824 & 0.00 & 0.00 & 0 & 0 & 0 \\
CIC-IDS2018 & K = 2 & 0.7824 & 0.00 & 0.00 & 0 & 0 & 0 \\
CIC-IDS2018 & K = 4 & 0.7824 & 0.00 & 0.00 & 0 & 0 & 0 \\
CIC-IDS2018 & K = 6 & 0.7824 & 0.00 & 0.00 & 0 & 0 & 0 \\
CIC-IDS2018 & K = 8 & 0.7824 & 0.00 & 0.00 & 0 & 0 & 0 \\
ToN-IoT & Global & 0.6796 & 0.00 & 0.00 & 0 & 0 & 0 \\
ToN-IoT & K = 2 & 0.7636 & 100.00 & 5.26 & 17,201 & 10,443 & 5,784 \\
ToN-IoT & K = 4 & 0.7617 & 100.00 & 1.23 & 6,504 & 5,381 & 1,101 \\
ToN-IoT & K = 6 & 0.7677 & 100.00 & 4.09 & 18,447 & 15,289 & 2,771 \\
ToN-IoT & K = 8 & 0.7463 & 100.00 & 0.69 & 12,290 & 8,328 & 3,741 \\
\bottomrule
\end{tabular}
\end{table}

\clearpage
\section{Research decisions and evidence status}
\begin{itemize}
\item The data-audit phase removed feature-identical conflicting-label rows for the current clean-data comparisons. Earlier datasets and results remain traceable in their original run folders.
\item EXP61 established clean-data baselines and measured cost. EXP63 supplied direct expert and scorer/verifier diagnostics over K=2,4,6,8. K=4 remains the previously used representative setting; no test-maximum K is substituted into the headline comparison.
\item On 2 October 2026, the presentation was revised around scarce-label attack detection. Aggregate F1 rankings are supporting evidence rather than the research objective.
\item The user clarified that existing Claude-generated scenario experiments were not the intended design. EXP64--68 are preserved as separate exploratory records; their existence does not imply adoption into this claim. EXP68 in particular is not the accepted sequential protocol.
\item Review of EXP68 found class-aware label selection within a window and evaluation using that window's collected labels. That procedure cannot establish strict predict-before-label streaming performance. Its field named benign\_fpr counts attacks predicted benign rather than false alarms on benign traffic. Original outputs are retained without silently rewriting their meaning.
\item Actual raw timestamps and few-label before/after adaptation comparisons are not mandatory. OOD means out-of-distribution data or new attacks and remains a later extension.
\end{itemize}
\paragraph{Record structure.} Raw runs remain in place. Research-direction versions provide relative-link views and metadata for protocol, validity and adoption status. This manuscript contains selected evidence; detailed HTML/Markdown records preserve diagnostic history. Generated tables and their source hashes are maintained in \texttt{sources.json}. Narrative sections are edited independently and are not overwritten when tables are regenerated.

\end{document}
